\documentclass[12pt]{article}
\usepackage{geometry}
\geometry{margin=1.4in}
\usepackage{setspace}
\setstretch{1.4}

\title{Learning as Displacement:\\
Why Schools Optimize for Grades, Not Growth}
\author{Diego Rincón\\
\small Independent}
\date{}

\begin{document}

\maketitle

\begin{abstract}
Education systems have an intended ground state: curious, capable, autonomous learners. They have a measurable proxy: grades. Schools displace from the first to the second because grades are cheaper to optimize for. The system works perfectly—at maximizing grades. It simply returns to the wrong ground state.
\end{abstract}

\section{The Learning Ground State}

True learning ground state $S_0$: students develop mastery, curiosity, metacognition, resilience, autonomy.

Measurable ground state $S^*$: students pass tests, get high grades, meet benchmarks.

Both require effort. The second is cheaper to verify.

\section{Why Schools Displace}

A teacher has limited time. They can:
- Teach for understanding (slow, requires ongoing assessment)
- Teach to the test (fast, results are measurable, tied to job security)

The system incentivizes the second. Not through malice, but through structure: funding is tied to test scores, teacher evaluations are tied to student grades, students' futures are tied to transcripts.

The cost to maintain $S_0$ is prohibitive. The system returns to $S^*$.

\section{The Student as Optimizer}

Once displaced, students optimize for grades:
- Memorizing instead of understanding
- Short-term cramming instead of long-term retention
- Test-taking strategy instead of knowledge
- Pleasing teachers instead of pursuing curiosity

This is not student failure. This is students finding the stable state and optimizing within it.

\section{The Irreversibility}

Once a student learns to optimize for grades, returning to curiosity-driven learning is expensive. They have built habits, expectations, and identities around grades. The system has selected for test-takers, not thinkers.

Breaking this requires:
- Removing grades (politically impossible)
- Changing incentives (systemically difficult)
- Individual heroics (a few teachers working against the system)

Most systems stay displaced.

\end{document}
